\section*{अध्याय \ref{ch:g12:adaptive-immunity} --- अनुकूली प्रतिरक्षा और टीकाकरण}

\begin{solution}{exo:g12:adaptive-immunity:1}
\omterm{def:g12:adaptive-immunity:antigen}{प्रतिजन} कोई भी ऐसा अणु है जिसे अनुकूली तंत्र पहचान सके; और \omterm{def:g12:adaptive-immunity:antigen}{लसीकाणु} ऐसी श्वेत
\omterm{def:g10:cells-common-unit:cell}{कोशिका} है जो किसी एक \omterm{def:g12:adaptive-immunity:antigen}{प्रतिजन} के लिए ग्राही ढोती है। उसकी विशिष्टता उसके
ग्राही से आती है, जो परिपक्वन के दौरान यादृच्छिक ढंग से जोड़ा जाता है और
उसकी सारी प्रतियों पर एक-सा होता है।
\end{solution}

\begin{solution}{exo:g12:adaptive-immunity:2}
वह \omterm{def:g12:adaptive-immunity:antigen}{प्रतिजन} उन थोड़े-से \omterm{def:g12:adaptive-immunity:antigen}{लसीकाणुओं} से बँधता है जिनका ग्राही बैठता हो; वे
\omterm{def:g10:cells-common-unit:cell}{कोशिकाएँ} दिनों पर किसी क्लोन में कई गुना होती हैं; वह क्लोन प्रभावक
\omterm{def:g10:cells-common-unit:cell}{कोशिकाओं} में विभेदित होता है जो लड़ती हैं; और एक अल्पसंख्या \omterm{prop:g12:adaptive-immunity:selection}{स्मृति कोशिकाएँ}
बनती है जो टिकी रहती हैं।
\end{solution}

\begin{solution}{exo:g12:adaptive-immunity:3}
प्लाज़्मा \omterm{def:g10:cells-common-unit:cell}{कोशिकाओं} का स्रावित किया कोई Y आकार का \omterm{def:g11:gene-expression:protein}{प्रोटीन}, यानी उस B \omterm{def:g10:cells-common-unit:cell}{कोशिका}
के ग्राही की घुलनशील प्रति; उसकी हर बाँह उस \omterm{def:g12:adaptive-immunity:antigen}{प्रतिजन} से बँधती है। किसी
जीवाणु पर वह ऐसी परत बनाता है जिसे \omterm{def:g12:innate-immunity:phagocytosis}{भक्षककोशिकाएँ} पकड़ती हैं, और जीवाणुओं को
आड़े-तिरछे जोड़कर ऐसे संकुल बनाता है जो निगल लिए जाते हैं।
\end{solution}

\begin{solution}{exo:g12:adaptive-immunity:4}
\omterm{prop:g12:adaptive-immunity:tcells}{सहायक T कोशिकाएँ}, जो \omterm{prop:g12:innate-immunity:handover}{वृक्षाभ कोशिकाओं} से सक्रिय होती हैं, वे साइटोकाइन
स्रावित करती हैं जो B \omterm{def:g10:cells-common-unit:cell}{कोशिकाओं} तथा बाक़ी T \omterm{def:g10:cells-common-unit:cell}{कोशिकाओं} को इजाज़त देते और बढ़ाते
हैं। और कोशिकाविषी T \omterm{def:g10:cells-common-unit:cell}{कोशिकाएँ} उन \omterm{def:g10:cells-common-unit:cell}{कोशिकाओं} को मारती हैं जो किसी विषाणु या
किसी \omterm{def:g11:cancer:cancer}{अर्बुद} \omterm{def:g11:gene-expression:protein}{प्रोटीन} के टुकड़े दिखा रही हों।
\end{solution}

\begin{solution}{exo:g12:adaptive-immunity:5}
उस रोगजनक के \omterm{def:g12:adaptive-immunity:antigen}{प्रतिजन} का कोई हानिरहित रूप (कमज़ोर किया, मारा हुआ, कोई
\omterm{def:g11:gene-expression:protein}{प्रोटीन}, कोई निष्क्रिय विष, या उसके निर्देश) किसी सहायक द्रव्य के साथ। वह
कोई प्राथमिक प्रतिक्रिया छेड़ता है और \omterm{prop:g12:adaptive-immunity:selection}{स्मृति कोशिकाएँ} छोड़ जाता है, ताकि वह
असली रोगजनक किसी द्वितीयक प्रतिक्रिया से मिले।
\end{solution}

\begin{solution}{exo:g12:adaptive-immunity:6}
प्राथमिक: एक हफ़्ते की देरी, दो से तीन हफ़्ते पर एक नीची चोटी (द्वितीयक की
लगभग दसवाँ हिस्सा), और महीने भर के भीतर गिरावट। द्वितीयक: दो-तीन दिन की
देरी, दस दिनों के भीतर दस गुनी ऊँची चोटी, जो महीनों टिकी रहती है।
\end{solution}

\begin{solution}{exo:g12:adaptive-immunity:7}
स्मृति विशिष्ट है: A की पहली ख़ुराक की छोड़ी \omterm{prop:g12:adaptive-immunity:selection}{स्मृति कोशिकाएँ} केवल A को जवाब
देती हैं; जबकि B, जो पहली बार मिला, कोई प्राथमिक प्रतिक्रिया पाता है।
\end{solution}

\begin{solution}{exo:g12:adaptive-immunity:8}
\omterm{prop:g12:adaptive-immunity:antibodies}{प्रतिरक्षी} तरलों में बैठे \omterm{def:g11:gene-expression:protein}{प्रोटीन} हैं और \omterm{def:g10:cells-common-unit:cell}{कोशिकाओं} में घुस नहीं सकते। कोई
संक्रमित \omterm{def:g10:cells-common-unit:cell}{कोशिका} उस विषाणु के टुकड़े अपनी सतह पर दिखाती है, और कोशिकाविषी T
\omterm{def:g10:cells-common-unit:cell}{कोशिकाएँ} उन्हें पहचानकर उस \omterm{def:g10:cells-common-unit:cell}{कोशिका} को मार देती हैं।
\end{solution}

\begin{solution}{exo:g12:adaptive-immunity:9}
$1 - 1/4 = 75\%$; $1 - 1/18 \approx 94\%$।
\end{solution}

\begin{solution}{exo:g12:adaptive-immunity:10}
उस दहलीज़ के ऊपर हर मामला एक से भी कम लोगों को संक्रमित करता है, इसलिए वह
विषाणु घूम ही नहीं पाता और शायद ही किसी शिशु तक पहुँचता है; प्रतिरक्षित
वयस्क ही उनके चारों ओर की दीवार हैं।
\end{solution}

\begin{solution}{exo:g12:adaptive-immunity:11}
प्रति माइक्रोलीटर लगभग 750, 520 और 180; और \omterm{prop:g12:adaptive-immunity:hiv}{एड्स} साल 9 के आसपास शुरू होता है,
जब वह गिनती 200 पार करती है।
\end{solution}

\begin{solution}{exo:g12:adaptive-immunity:12}
\omterm{prop:g12:adaptive-immunity:tcells}{सहायक T कोशिकाओं} के साइटोकाइन ही वे हैं जो B \omterm{def:g10:cells-common-unit:cell}{कोशिकाओं} को प्लाज़्मा \omterm{def:g10:cells-common-unit:cell}{कोशिकाएँ}
बनने की और कोशिकाविषी T \omterm{def:g10:cells-common-unit:cell}{कोशिकाओं} को कई गुना होने की इजाज़त देते हैं; दोनों
बाँहें उसी एक समन्वयक पर निर्भर हैं, इसलिए उसे हटाना दोनों को चुप करा देता है।
\end{solution}

\begin{solution}{exo:g12:adaptive-immunity:13}
वे \omterm{prop:g12:adaptive-immunity:antibodies}{प्रतिरक्षी} उस विषाणु के ख़िलाफ़ चुने गए थे जैसा वह था; और वह विषाणु अपने
सतह \omterm{def:g11:gene-expression:protein}{प्रोटीन} नए क्लोनों के चुने जाने से तेज़ उत्परिवर्तित करता है, इसलिए हर
प्रतिक्रिया पीछे छूट जाती है। और वह विषाणु ठीक उन्हीं सहायक \omterm{def:g10:cells-common-unit:cell}{कोशिकाओं} के भीतर
छिपता तथा कई गुना होता है जो उस प्रतिक्रिया का समन्वय करतीं।
\end{solution}

\begin{solution}{exo:g12:adaptive-immunity:14}
\omterm{prop:g12:adaptive-immunity:selection}{स्मृति कोशिकाएँ} वही सतह \omterm{def:g11:gene-expression:protein}{प्रोटीन} पहचानती हैं जो उस \omterm{def:g12:adaptive-immunity:vaccine}{टीके} ने पेश किए थे।
इन्फ़्लुएंज़ा के हर साल ऐसे रूपों में उत्परिवर्तित हो जाते हैं जिनसे पुरानी
स्मृति बँधती नहीं; जबकि ख़सरे के सतह \omterm{def:g11:gene-expression:protein}{प्रोटीन} शायद ही बदलते हैं, इसलिए 1970
की स्मृति आज के विषाणु पर भी बैठती है।
\end{solution}

\begin{solution}{exo:g12:adaptive-immunity:15}
क्या बदलता है: \omterm{def:g12:adaptive-immunity:antigen}{लसीकाणुओं} के ग्राही, जो यादृच्छिक ढंग से बने हैं। क्या चुनता
है: वह \omterm{def:g12:adaptive-immunity:antigen}{प्रतिजन}, जो केवल बैठते क्लोनों को कई गुना होने देता है। क्या विरासत
में मिलता है: वह ग्राही, उस क्लोन के वंशजों को, जिनमें \omterm{prop:g12:adaptive-immunity:selection}{स्मृति कोशिकाएँ} भी
हैं। समय-पैमाना: किसी प्रतिक्रिया के लिए दिन, स्मृति के लिए एक जीवन, और
इसके सामने किसी \omterm{def:g12:selection-drift-speciation:population}{आबादी} के लिए पीढ़ियाँ। वह उपमा संचरण पर रुक जाती है: चुने
हुए क्लोन उस शरीर के साथ मर जाते हैं और अगली पीढ़ी तक नहीं जाते।
\end{solution}

\begin{solution}{pb:g12:adaptive-immunity:1}
\textbf{1.} लगभग एक हफ़्ता कुछ नहीं, दो से तीन हफ़्ते पर एक मामूली चोटी, और
आगे के हफ़्तों पर गिरावट।

\textbf{2.} दो-तीन दिनों के भीतर \omterm{prop:g12:adaptive-immunity:antibodies}{प्रतिरक्षी}, दस गुनी ऊँची चोटी, जो महीनों
चलती है: पहली ख़ुराक की छोड़ी \omterm{prop:g12:adaptive-immunity:selection}{स्मृति कोशिकाएँ} पहले ही एक बड़ा क्लोन हैं,
इसलिए गुणन कुछ ही \omterm{def:g10:cells-common-unit:cell}{कोशिकाओं} के बजाय हज़ारों \omterm{def:g10:cells-common-unit:cell}{कोशिकाओं} से शुरू होता है।

\textbf{3.} कोई जीवित विषाणु \omterm{def:g10:cells-common-unit:cell}{कोशिकाओं} के भीतर थोड़ी देर कई गुना होता है,
इसलिए संक्रमित \omterm{def:g10:cells-common-unit:cell}{कोशिकाएँ} उसके टुकड़े दिखाती हैं और \omterm{prop:g12:adaptive-immunity:antibodies}{प्रतिरक्षियों} के साथ-साथ
कोशिकाविषी T \omterm{def:g10:cells-common-unit:cell}{कोशिकाएँ} भी चुनी जाती हैं; जबकि कोई मारा हुआ विषाणु मुख्यतः
\omterm{prop:g12:adaptive-immunity:antibodies}{प्रतिरक्षी} वाली बाँह सिखाता है।

\textbf{4.} पहले से मौजूद \omterm{prop:g12:adaptive-immunity:antibodies}{प्रतिरक्षी} उस विषाणु का अधिकांश रोक देते हैं;
स्मृति B तथा T \omterm{def:g10:cells-common-unit:cell}{कोशिकाएँ} दो दिनों के भीतर कई गुना हो जाती हैं; और संक्रमित
\omterm{def:g10:cells-common-unit:cell}{कोशिकाएँ} विषाणु के फैलने से पहले मार दी जाती हैं। वह संक्रमण लक्षणों से पहले
ही रोक दिया जाता है।

\textbf{5.} तीन दिन चुपचाप गुणन, फिर बुख़ार, ददोरे तथा बीमारी, जबकि
प्राथमिक प्रतिक्रिया हफ़्ते भर पर बनती है; और दिन दस के आसपास सेहतयाबी,
\omterm{prop:g12:adaptive-immunity:selection}{स्मृति कोशिकाओं} तथा जीवन भर की प्रतिरक्षा के साथ --- यदि वह उन जटिलताओं से
बच जाए।

\textbf{6.} $1 - 1/15 \approx 93\%$।

\textbf{7.} $15 \times 0.10 = 1.5$: यानी हर मामला एक से अधिक लोगों को
संक्रमित करता है, इसलिए वह बीमारी फैल सकती है। सुरक्षित नहीं।

\textbf{8.} $10\% \times \num{50000} = 5000$ प्रतिरक्षित नहीं, जिनमें से
1500 वे शिशु हैं।

\textbf{9.} 2500 ग्राह्य; $15 \times 0.05 = 0.75$: यानी एक से कम, और वह
बीमारी बुझ जाती है।

\textbf{10.} शिशुओं को \omterm{def:g12:adaptive-immunity:vaccine}{टीका} नहीं लगाया जा सकता, इसलिए उनकी इकलौती रक्षा यही
है कि उनके इर्द-गिर्द के वयस्क, प्रतिरक्षित होकर, उस विषाणु को उन तक
पहुँचने का कोई रास्ता ही न छोड़ें।

\textbf{11.} पीढ़ी 3: $1.5^3 \approx 3.4$ मामले; और पीढ़ी 6:
$1.5^6 \approx 11$।

\textbf{12.} $1 + 1.5 + 2.25 + 3.4 + 5.1 + 7.6 + 11.4 \approx 32$ मामले।

\textbf{13.} 32 का लगभग 30\%: यानी कोई 10 शिशु।

\textbf{14.} हर मामला एक ग्राह्य व्यक्ति हटाता है और एक प्रतिरक्षित जोड़ता
है, इसलिए ग्राह्य अंश गिरता है और उसके साथ हर मामले के संक्रमितों की संख्या
भी; 5000 ग्राह्य लोगों के साथ वह प्रकोप सैकड़ों मामलों के बाद ही, महीनों पर,
धीमा पड़ता --- बशर्ते टीकाकरण बीच में न आए।

\textbf{15.} 95\% पर: $0.75^3 \approx 0.4$ और $0.75^6 \approx 0.18$
मामले; कुल लगभग $1 + 0.75 + 0.56 + \dots \approx 4$ मामले। यानी दस गुना कम
मामले, और वह प्रकोप ख़ुद ही बुझ जाता है।

\textbf{16.} 8 सालों में 800: यानी प्रति माइक्रोलीटर प्रति साल 100, और प्रति
महीने लगभग 8।

\textbf{17.} B \omterm{def:g10:cells-common-unit:cell}{कोशिकाओं} को प्लाज़्मा \omterm{def:g10:cells-common-unit:cell}{कोशिकाएँ} बनने के लिए सहायक \omterm{def:g10:cells-common-unit:cell}{कोशिकाओं} के
साइटोकाइन चाहिए, और कोशिकाविषी T \omterm{def:g10:cells-common-unit:cell}{कोशिकाओं} को कई गुना होने के लिए भी: इसलिए
बहुत कम सहायकों के साथ किसी भी बाँह को इजाज़त नहीं मिलती, हालाँकि उनकी
\omterm{def:g10:cells-common-unit:cell}{कोशिकाएँ} मौजूद हैं।

\textbf{18.} हाँ: \omterm{prop:g12:adaptive-immunity:selection}{स्मृति कोशिकाओं} को किसी द्वितीयक प्रतिक्रिया में फैलने के
लिए आज भी सहायक साइटोकाइन चाहिए; और समन्वय के बिना वह स्मृति इस्तेमाल ही
नहीं हो सकती, तथा ख़सरा, जो किसी स्वस्थ वयस्क के लिए हानिरहित है, जानलेवा
हो सकता है।

\textbf{19.} सहायक \omterm{def:g12:selection-drift-speciation:population}{आबादी} बची \omterm{def:g10:cells-common-unit:cell}{कोशिकाओं} तथा मज्जा से दोबारा उग आती है, और
बाक़ी बाँहों की \omterm{prop:g12:adaptive-immunity:selection}{स्मृति कोशिकाएँ} कभी नष्ट हुई ही नहीं थीं --- यानी समन्वय
लौटते ही वह भंडार सही-सलामत है। पर वह विषाणु छिपी \omterm{def:g10:cells-common-unit:cell}{कोशिकाओं} में टिका रहता है,
इसलिए दवाएँ रोकने पर वह फिर चल पड़ता है।

\textbf{20.} 93\% कवरेज; 90\% पर किसी एक यात्री से लगभग 32 मामले; और वह
\omterm{prop:g12:adaptive-immunity:tcells}{सहायक T कोशिका}।
\end{solution}
